
\begin{exercise}[subtitle={Randomized linear algebra \Coffeecup}]
    Let \(x,y \in \real^\dims\), where the dimension \(\dims\) is really large.
    The computation of \(\scp{x,y}\) clearly costs \(\dims\) multiplications and
    additions. We will try to speed it up using a random sample.
    Specifically, let \(Z_1, \ldots Z_n\) be independently drawn from
    \(\frac1\dims\sum_{i=1}^\dims \delta_{x_iy_i}\), that is
    \(\prob{Z_1 = x_iy_i} = \frac1\dims\) for \(i=1, \ldots, \dims\).
    \begin{enumerate}[label=(\alph*), noitemsep]
        \item\label{it: (randomized linear algebra) expectation}
        Prove that \(\dims \E{Z_1} = \scp{x,y}\),

        \item\label{it: (randomized linear algebra) convergence}
        Assuming \(\scp{x,y} \neq 0\), prove the following estimator
        \[
            \hat S_n \coloneq \frac\dims{n} \sum_{i=1}^n Z_i,
        \]
        converges against the scalar product in relative terms, that is
        \[
            \frac{\hat S_n}{\scp{x,y}} \overset\as\to 1 \qquad (n \to \infty).
        \]

        \item\label{it: (randomized linear algebra) PAC bound}
        For \(\hat S_n\) to be an effective speedup of the
        computation of the scalar product it is necessary that \(n\le \dims\). So
        clearly \(n\to \infty\) is not an interesting statement.
        Let \(m = \max_i \abs{x_i y_i}\) and show that
        \[
            \prob[\Big]{\abs[\Big]{\frac{\hat S_n}{\scp{x,y}} - 1} > \epsilon}
            \le 2\exp\Bigl(-n \epsilon^2 \frac{\E{Z_1}^2}{8 m^2}\Bigr).
        \]
        \textbf{Hint:} Corollary \ref{cor: hoeffding inequality for bounded random variables}.
        \begin{remark}[Many significant values]
            If most \(x_iy_i\) are very small except for a few \(i=\set{1,\dots, \dims}\),
            then the inner product \(\scp{x,y}\) is dominated by these values.
            This means that a random sample can be very far off if it does not
            contain these significant indices. This case is captured
            by the ratio \(\frac{\E{Z_1}}{m}\): If the average product
            \(\E{Z_1}\) is much smaller, than largest product \(m\), the sum is
            dominated by a few large entries and convergence is slow.
        \end{remark}

        \item\label{it: (randomized linear algebra) sample complexity}
        Assume \(\frac{\E{Z_1}}{m} \ge 0.1\), that a precision of \(\epsilon
        = 10^{-4}\) is required and a \(99\%\) confidence. How big must
        \(\dims\) be such that this approach reduces the computational cost?
        How big to get a \(99.999\%\) confidence?
    \end{enumerate}
\end{exercise}

\begin{solution}
\begin{enumerate}[wide=0pt]
    \item[\ref{it: (randomized linear algebra) expectation}]
    By definition we have
    \[
        \E{Z_1} = \frac1\dims\sum_{i=1}^\dims  x_i y_i = \frac1\dims \scp{x,y}.
    \]            

    \item[\ref{it: (randomized linear algebra) convergence}]
    We have \(\frac{\hat S_n}{\scp{x,y}} = \frac1n \sum_{i=1}^n Z_i /
    \E{Z_1}\) and thus the claim by the strong law of large numbers.

    \item[\ref{it: (randomized linear algebra) PAC bound}]
    Since \(Z_i \in [-m,m]\) for \(i=1, \ldots, n\) we can apply
    Hoeffding's inequality for bounded random variables (Corollary \ref{cor:
    hoeffding inequality for bounded random variables}) to get
    \begin{align*}
        \prob[\Big]{\abs[\Big]{\frac{\hat S_n}{\scp{x,y}} - 1} > \epsilon}
        &=\prob[\Big]{\abs[\Big]{\frac1n\sum_{i=1}^n Z_i - \E{Z_1}} > \epsilon \abs{\E{Z_1}}}
        \\
        &\le 2\exp\Bigl(-n \epsilon^2 \frac{\E{Z_1}^2}{2(2m)^2}\Bigr).
    \end{align*}

    \item[\ref{it: (randomized linear algebra) sample complexity}]
    We need \(n\) to satisfy
    \[
        2\exp\Bigl(-n \epsilon^2 \frac{\E{X_1 Y_1}^2}{8 m^2}\Bigr) \le \delta 
        \iff n \ge \frac{8\ln(\frac1{2\delta})}{\epsilon^2 \frac{\E{X_1 Y_1}^2}{m^2}}.
    \]
    Since \(\E{X_1 Y_1}^2 \ge 0.1\) we have that
    \[
        n \ge \frac{80\ln(\frac1{2\delta})}{\epsilon^2}
    \]
    is sufficient to get the required confidence. We now simply
    plug in the numbers. For \(\delta = 0.01 = 1\%\) we have
    \[
        n \ge \frac{80\ln(\frac1{0.02})}{10^{-8}} \approx 3.12 \cdot 10^{10}.
    \]
    For \(\delta = 0.00001 = 0.001\%\) we have
    \[
        n \ge \frac{80\ln(\frac1{0.00002})}{10^{-8}} \approx 8.66 \cdot 10^{10}.
    \]
    If \(n\le \dims\) then this approach reduces the computational cost.
    So for \(\dims \ge 3.12 \cdot 10^{10}\) (or \(\dims \ge 8.66 \cdot 10^{10}\) respectively)
    this approach reduces the computational cost. As in Remark \ref{rem:
    accuracy vs confidence} the confidence has almost no effect on the
    required dimension. In contrast the precision \(\epsilon\)
    dominates the required dimension.
    \qedhere
\end{enumerate}
\end{solution}
